\documentclass[activite]{thedocument}
\usepackage{enumitem}
\startdatakeys
\source{originale}
\BO{2026}
\cycle{4}
\competencesDuSocle{RA3}
\objectifsApprentissage{C4B3ca}
\motsCles{réciproque du théorème de Thalès, parallélisme}
\enddatakeys
\begin{document}
    En observant attentivement la 
    \videoManim{reciproque_thales}{DiscoverReciproqueThales}{vidéo projetée au tableau},
    répondre aux questions suivantes.\par\noindent
    {\centering\Large\textbf{Partie I}\par}
    \bigskip\noindent
    À l'\oe il nu, les droites sont-elles parallèles~?
    \begin{enumerate}[label=Q\arabic* :, font=\bfseries]
        \item
        \item
        \item
        \item
        \item
    \end{enumerate}
    {\centering\Large\textbf{Partie II}\par}
    \bigskip\noindent
    À l'\oe il nu, les droites sont-elles parallèles~? Attention, la droite
    \mathspoint{(BD)} est cachée.
    \begin{enumerate}[label=Q\arabic* :, font=\bfseries]
        \item
        \item
        \item
        \item
        \item
    \end{enumerate}
    {\centering\Large\textbf{Partie III}\par}
    \bigskip\noindent
    Par le calcul, les droites sont-elles parallèles~?\par\noindent
    \answerspace{5 cm}
\end{document}